\documentclass[11pt,a4paper]{article}
\usepackage[a4paper,margin=2.5cm]{geometry}
\usepackage{amsmath,amssymb,amsfonts,mathtools}
\usepackage{bm}
\usepackage{booktabs}
\usepackage{enumitem}
\usepackage{hyperref}
\hypersetup{colorlinks=true,linkcolor=blue,citecolor=blue,urlcolor=blue}

\newcommand{\W}{\bm W}
\newcommand{\A}{\bm A}
\newcommand{\U}{\bm U}
\newcommand{\K}{\bm K}
\newcommand{\D}{\bm D}

\title{Gu\'ia Did\'actica del M\'etodo MAW--ECM\\
\large Basada en \texttt{DOCS\_MAW\_ECM} (enfoque profesor)}
\author{Notas de estudio para implementaci\'on}
\date{\today}

\begin{document}
\maketitle

\begin{abstract}
Este documento explica, de forma did\'actica y paso a paso, la l\'ogica del m\'etodo
\textbf{MAW--ECM} tal como aparece en la carpeta \texttt{DOCS\_MAW\_ECM}:
\texttt{MAW\_ECM\_genV1\_DOC.tex}, \texttt{MAW\_ECM\_genV1.m},
\texttt{GraphBased\_MAW\_ECM\_2v.m}, \texttt{prune\_step\_optionB.m},
\texttt{MAWecmBasisMATRIX\_LOCAL\_noverlap.m} y
\texttt{RegressionViaRBF\_master\_weights.m}.

El objetivo es entender claramente:
\begin{enumerate}[leftmargin=1.2cm]
    \item qu\'e problema matem\'atico resuelve,
    \item por qu\'e se usa \emph{pruning} iterativo y no una soluci\'on global directa,
    \item c\'omo entran la positividad, los active sets y el Laplaciano de grafo,
    \item y c\'omo se obtiene finalmente $w(\bm q_M)$ para uso online.
\end{enumerate}
\end{abstract}

\tableofcontents

\section{Diccionario completo de s\'imbolos (muy expl\'icito)}
\subsection{\'Indices y tama\~nos}
\begin{itemize}[leftmargin=0.8cm]
    \item $i$: \'indice de \textbf{elemento} (fila de pesos).
    \item $j$: \'indice de muestra del manifold latente (columna de pesos).
    \item $n_{\text{ip}}$: en el documento original, n\'umero de puntos de integraci\'on; en nuestra lectura pr\'actica, \textbf{n\'umero de elementos candidatos} del ensamblaje hiperreducido.
    \item $N_M$: n\'umero de muestras latentes offline.
    \item $n_c$: n\'umero de elementos candidatos activos en una etapa de pruning.
    \item $m_j$: n\'umero de restricciones locales en la muestra $j$ (dimensi\'on de $\bm b_j$).
\end{itemize}

\subsection{Pesos: cu\'ales hay y qu\'e significan}
\begin{itemize}[leftmargin=0.8cm]
    \item $\bm w^{\text{FE}}\in\mathbb R^{n_{\text{ip}}}$:
    pesos base asociados a cada elemento candidato (en formulaci\'on cl\'asica se interpretan como pesos de cuadratura).
    \item $\bm w_0$ o $\bm w_{\text{INI}}$:
    pesos iniciales de ECM fijo sobre el soporte inicial $\mathcal Z_0$.
    \item $\bm w^{(j)}$:
    pesos en la muestra latente $j$.
    \item $\W\in\mathbb R^{n_c\times N_M}$:
    matriz de pesos adaptativos; fila $i$ = evoluci\'on del peso del punto $i$ en todo el manifold; columna $j$ = vector de pesos en el estado $j$.
\end{itemize}

\subsection{Matrices de integrandos/restricciones}
\begin{itemize}[leftmargin=0.8cm]
    \item $\A(\bm q_M)\in\mathbb R^{n_{\text{ip}}\times m}$:
    matriz de \textbf{contribuciones elementales}; cada fila corresponde a un elemento y cada columna a una cantidad reducida que quieres reproducir (en nuestro caso, componentes del residual proyectado).
    \item $\U_j\in\mathbb R^{n_c\times m_j}$:
    operador local restringido al soporte candidato en la muestra $j$.
    \item $\bm b_j\in\mathbb R^{m_j}$:
    vector objetivo local (la suma ``full'' que queremos reproducir con soporte reducido) en la muestra $j$.
\end{itemize}

\subsection{La duda clave: qu\'e es $A_j$ y por qu\'e aparece $U_j^T$}
En los documentos del profesor aparecen \textbf{dos notaciones equivalentes}:
\[
\U_j^T \bm w^{(j)} = \bm b_j
\quad\Longleftrightarrow\quad
\A_j \bm w^{(j)} = \bm b_j
\]
con
\[
\A_j := \U_j^T.
\]

Esto pasa porque en c\'odigo MATLAB de la etapa de pruning se construye
expl\'icitamente:
\[
\texttt{A\{j\} = U\{j\}(...,: )'}.
\]

Entonces:
\begin{itemize}[leftmargin=0.8cm]
    \item si usas $\U_j$, la ecuaci\'on es $\U_j^T\bm w^{(j)}=\bm b_j$;
    \item si renombras $\A_j=\U_j^T$, la misma ecuaci\'on se escribe $\A_j\bm w^{(j)}=\bm b_j$.
\end{itemize}

\subsection{Chequeo de dimensiones (para no perderse)}
Si
\[
\U_j\in\mathbb R^{n_c\times m_j},\qquad
\bm w^{(j)}\in\mathbb R^{n_c},\qquad
\bm b_j\in\mathbb R^{m_j},
\]
entonces
\[
\U_j^T\in\mathbb R^{m_j\times n_c},
\]
y por tanto
\[
\U_j^T\bm w^{(j)}\in\mathbb R^{m_j},
\]
que coincide con el tama\~no de $\bm b_j$.

\subsection{Traducci\'on f\'isica de la ecuaci\'on local}
\[
\U_j^T\bm w^{(j)}=\bm b_j
\]
significa: ``en el estado latente $j$, los pesos elegidos deben reproducir
exactamente (o casi exactamente) los $m_j$ valores objetivo de \textbf{residual reducido}''.

\subsection{Transpuesta: recordatorio corto}
Para una matriz
\[
M=
\begin{bmatrix}
a & b & c\\
d & e & f
\end{bmatrix}
\in\mathbb R^{2\times 3},
\]
su transpuesta es
\[
M^T=
\begin{bmatrix}
a & d\\
b & e\\
c & f
\end{bmatrix}
\in\mathbb R^{3\times 2}.
\]
Se intercambian filas por columnas.

\subsection{Mini-ejemplo num\'erico de $U_j^T w = b_j$}
Sup\'on $n_c=4$, $m_j=2$:
\[
\U_j=
\begin{bmatrix}
1 & 0\\
0 & 1\\
1 & 1\\
2 & 1
\end{bmatrix},
\qquad
\bm w^{(j)}=
\begin{bmatrix}
0.2\\0.1\\0.3\\0.4
\end{bmatrix}.
\]
Entonces
\[
\U_j^T\bm w^{(j)}=
\begin{bmatrix}
1 & 0 & 1 & 2\\
0 & 1 & 1 & 1
\end{bmatrix}
\begin{bmatrix}
0.2\\0.1\\0.3\\0.4
\end{bmatrix}
=
\begin{bmatrix}
1.3\\0.8
\end{bmatrix}
=\bm b_j.
\]
Si ese $\bm b_j$ era el objetivo local, la muestra $j$ es factible.

\subsection{Otros s\'imbolos que ver\'as todo el tiempo}
\begin{itemize}[leftmargin=0.8cm]
    \item $\K_G$: Laplaciano del grafo en espacio latente (penaliza cambios bruscos entre vecinos).
    \item $\alpha_G$: intensidad de regularizaci\'on de grafo.
    \item $\Delta\W=\W_{\text{new}}-\W_{\text{old}}$: cambio de pesos en una poda tentativa.
    \item $\mathcal D_j$: conjunto activo local (\'indices clampeados a cero en nodo $j$).
    \item $\mathcal Z_0,\mathcal Z_{\text{red}}$: soporte inicial e integrado final.
\end{itemize}

\section{Idea central en una frase}
MAW--ECM mantiene un \textbf{soporte com\'un reducido} de elementos,
pero permite que los \textbf{pesos cambien con la coordenada latente}:
\[
\W(\bm q_M)=
\begin{bmatrix}
w_1(\bm q_M)\\
\vdots\\
w_{n_{\text{red}}}(\bm q_M)
\end{bmatrix}.
\]

En lugar de un ECM cl\'asico con pesos fijos, aqu\'i hay un campo de pesos sobre la
variedad (manifold) reducida.

\section{Recordatorio de ECM cl\'asico}
Sea $\mathcal Z_{\text{FE}}=\{1,\dots,n_{\text{ip}}\}$ el conjunto completo de
elementos candidatos del modelo completo, y $\bm w^{\text{FE}}$ su vector de pesos.
Para un conjunto de contribuciones reducidas en una matriz $\A(\bm q_M)\in\mathbb R^{n_{\text{ip}}\times m}$:
\[
\bm b(\bm q_M)=\A(\bm q_M)^T\bm w^{\text{FE}}.
\]

Aqu\'i, dicho sin abreviaturas:
\begin{itemize}[leftmargin=0.8cm]
    \item $n_{\text{ip}}$ = n\'umero total de elementos candidatos.
    \item $m$ = n\'umero de cantidades reducidas que quieres reproducir (por ejemplo, componentes del residual proyectado).
    \item $\A(\bm q_M)$ tiene \textbf{filas} asociadas a elementos y \textbf{columnas} asociadas a esas $m$ cantidades.
    \item $\bm w^{\text{FE}}$ son los pesos FE ``de toda la vida''.
    \item $\bm b(\bm q_M)$ es el vector objetivo que define lo que el modelo reducido debe reproducir en ese estado latente $\bm q_M$.
\end{itemize}

Ejemplo de relleno (s\'olo did\'actico): si $n_{\text{ip}}=4$ y $m=2$,
\[
\A(\bm q_M)=
\begin{bmatrix}
a_{11} & a_{12}\\
a_{21} & a_{22}\\
a_{31} & a_{32}\\
a_{41} & a_{42}
\end{bmatrix},
\qquad
\bm w^{\text{FE}}=
\begin{bmatrix}
w_1^{\text{FE}}\\ w_2^{\text{FE}}\\ w_3^{\text{FE}}\\ w_4^{\text{FE}}
\end{bmatrix}.
\]
Entonces
\[
\A(\bm q_M)^T\bm w^{\text{FE}}=
\begin{bmatrix}
a_{11}w_1^{\text{FE}}+a_{21}w_2^{\text{FE}}+a_{31}w_3^{\text{FE}}+a_{41}w_4^{\text{FE}}\\
a_{12}w_1^{\text{FE}}+a_{22}w_2^{\text{FE}}+a_{32}w_3^{\text{FE}}+a_{42}w_4^{\text{FE}}
\end{bmatrix}.
\]

ECM cl\'asico busca un soporte fijo $\mathcal Z_{\text{red}}$ y pesos fijos
$\bm w_{\text{red}}$ tales que:
\[
\A_{\mathcal Z_{\text{red}}}(\bm q_M)^T \bm w_{\text{red}}
\approx
\A(\bm q_M)^T\bm w^{\text{FE}}
\quad \forall \bm q_M \text{ de entrenamiento}.
\]

Problema: un solo vector $\bm w_{\text{red}}$ debe servir para todos los estados;
eso suele forzar soportes m\'as grandes.

\section{Problema ideal de MAW--ECM (versi\'on pr\'actica)}
Para cada muestra latente $\bm q_M^{(j)}$, $j=1,\dots,N_M$, se define un sistema local:
\[
\U_j^T \bm\zeta^{(j)}=\bm b_j,\qquad \bm\zeta^{(j)}\ge 0.
\]

Traducci\'on literal:
\begin{itemize}[leftmargin=0.8cm]
    \item $j$ = muestra latente (por ejemplo, un punto de la malla $q_{M1}\times q_{M2}$).
    \item $\bm\zeta^{(j)}$ = pesos de cubatura para esa muestra $j$.
    \item $\U_j$ = matriz local de contribuciones elementales del sistema reducido (por ejemplo residual proyectado) en el estado $j$.
    \item $\bm b_j$ = vector objetivo local de ese mismo sistema en $j$.
    \item $\U_j^T\bm\zeta^{(j)}=\bm b_j$ = exactitud local.
\end{itemize}

En varios scripts del profesor la misma ecuaci\'on aparece como
\[
\A_j\bm\zeta^{(j)}=\bm b_j,\qquad \text{con } \A_j:=\U_j^T.
\]
No son dos cosas distintas: es la \textbf{misma ecuaci\'on con distinto nombre}.

Si $\bm\zeta^{(j)}\in\mathbb R^{n_{\text{ip}}}$ son los pesos por elemento en el estado $j$,
el objetivo conceptual (soporte com\'un + suavidad) es:
\[
\min_{\{\bm\zeta^{(j)}\}}
\left\|\sum_{j=1}^{N_M}\bm\zeta^{(j)}\right\|_0
\;+\;
\frac{\alpha_G}{2}\sum_{i=1}^{n_{\text{ip}}}\bm\zeta_i^T\K_G\bm\zeta_i
\]
sujeto a exactitud local y positividad.

\textbf{Lectura f\'isica}:
\begin{itemize}[leftmargin=0.8cm]
    \item $\|\cdot\|_0$: minimizar la uni\'on de soportes (pocos elementos globales).
    \item $\K_G$: penalizar que los pesos ``salten'' entre vecinos del manifold.
    \item restricciones: conservar exactitud local del residual reducido y admisibilidad ($w\ge 0$).
\end{itemize}

Este problema global es combinatorio + restringido + acoplado: muy costoso para
resolver ``de una''.

\section{Estrategia real del profesor: pruning factible}
El flujo implementado no resuelve el problema global directamente.
Hace esto:
\begin{enumerate}[leftmargin=1.2cm]
    \item parte de una regla ECM fija factible inicial $(\mathcal Z_0,\bm w_0)$;
    \item replica $\bm w_0$ en todas las muestras del manifold;
    \item intenta eliminar un punto por vez (\emph{single-point pruning});
    \item redistribuye pesos para mantener exactitud (y luego positividad);
    \item si factible, acepta la eliminaci\'on; si no, la rechaza.
\end{enumerate}

Es exactamente el procedimiento operativo visible en \texttt{MAW\_ECM\_genV1.m}.

\subsection{Construcci\'on de sistemas locales}
En \texttt{MAWecmBasisMATRIX\_LOCAL\_noverlap.m} se construye, para cada muestra
$j$:
\[
\U_j,\;\bm b_j,\;\bm w^{(j)}_{\text{init}}.
\]

La idea pr\'actica es que $\U_j^T\bm w^{(j)}=\bm b_j$ represente
``lo que el residual proyectado debe reproducir'' en ese estado del manifold.

Aqu\'i tambi\'en es importante distinguir:
\begin{itemize}[leftmargin=0.8cm]
    \item $\bm w^{(j)}_{\text{init}}$: pesos iniciales (copiados del ECM fijo).
    \item $\bm w^{(j)}$: pesos que se ir\'an adaptando cuando empiece el pruning.
    \item Si en c\'odigo ves \texttt{U\{j\}(ilocPOS,:)'}, ese ap\'ostrofe es transpuesta:
    de ah\'i nace $A_j=\U_j^T$.
\end{itemize}

\subsection{Matriz de pesos adaptativos}
Si hay $n_c=|\mathcal Z_0|$ candidatos y $N_M$ muestras:
\[
\W \in \mathbb R^{n_c\times N_M},\qquad
\W=
\begin{bmatrix}
| & | & & |\\
\bm w^{(1)} & \bm w^{(2)} & \cdots & \bm w^{(N_M)}\\
| & | & & |
\end{bmatrix}.
\]
Eliminar un elemento de soporte equivale a eliminar una fila global de $\W$.

\section{Etapa 1: redistribuci\'on sin imponer positividad}
Script base: \texttt{Loop\_MAWecmNOENF\_cl.m}.

Para un candidato $p$ a eliminar, se impone $w_{p}^{(j)}=0$ y se resuelve, por
nodo/manifold $j$:
\[
\min_{\widetilde{\bm w}^{(j)}}
\frac{1}{2}\|\widetilde{\bm w}^{(j)}-\widetilde{\bm w}_{\text{old}}^{(j)}\|_2^2
\quad
\text{s.a.}\quad
\widetilde{\U}_j^T\widetilde{\bm w}^{(j)}=\bm b_j.
\]

El ``tilde'' significa ``despu\'es de quitar el candidato $p$'':
\begin{itemize}[leftmargin=0.8cm]
    \item $\widetilde{\bm w}^{(j)}$ = vector de pesos sin la entrada del punto eliminado.
    \item $\widetilde{\U}_j$ = la misma restricci\'on local, pero con la fila del punto eliminado fuera.
\end{itemize}

Soluci\'on de ``m\'inimo cambio'':
\[
\widetilde{\bm w}^{(j)}_{\text{new}}
=
\widetilde{\bm w}^{(j)}_{\text{old}}
-\widetilde{\U}_j\bm\lambda_j,
\quad
(\widetilde{\U}_j^T\widetilde{\U}_j)\bm\lambda_j
=
\widetilde{\U}_j^T\widetilde{\bm w}^{(j)}_{\text{old}}-\bm b_j.
\]

En c\'odigo, el incremento se obtiene con \texttt{lsqminnorm}.
Si aparece un peso negativo, la prueba se rechaza en esta etapa.

\section{Etapa 2: positividad local por active set}
Script base: \texttt{Loop\_MAWecmNOENFe.m} (l\'ogica descrita en el DOC).

Cuando la etapa 1 ya no puede podar, se impone:
\[
\widetilde{\bm w}^{(j)}\ge 0.
\]

Para cada nodo $j$ se usa un conjunto activo local $\mathcal D_j$:
\[
i\in\mathcal D_j \iff w_i^{(j)}=0 \text{ (clamp)}.
\]

Esto quiere decir:
\begin{itemize}[leftmargin=0.8cm]
    \item $\mathcal D_j$ no es global: cambia con $j$.
    \item un peso puede estar clampeado en un nodo latente y libre en otro.
    \item por eso $w(\bm q_M)$ suele quedar a trozos (piecewise), no perfectamente suave.
\end{itemize}

Se resuelve iterativamente:
\begin{enumerate}[leftmargin=1.2cm]
    \item resolver en variables libres $\mathcal I^{(j)}_{\text{free}}$;
    \item detectar componentes negativas;
    \item a\~nadirlas a $\mathcal D_j$;
    \item repetir hasta converger o fallar.
\end{enumerate}

\textbf{Punto clave}: el soporte sigue siendo com\'un, pero los active sets son
distintos por nodo. Eso explica no suavidad de $w(\bm q_M)$.

\section{Etapa 3: positividad con regularizaci\'on de grafo}
Scripts: \texttt{GraphBased\_MAW\_ECM\_2v.m} y
\texttt{prune\_step\_optionB.m}.

\subsection{Laplaciano sobre el espacio latente}
Para el caso 2D, se arma una malla cuadrilateral en $(q_{M1},q_{M2})$ y su
``stiffness'' escalar:
\[
\K_G \in \mathbb R^{N_M\times N_M},
\]
con scripts \texttt{quadMeshFromXY.m} y \texttt{ComputeK\_scl.m}.

Interpretaci\'on: $\K_G$ juega el rol de Laplaciano discreto de grafo.
Si dos muestras latentes son vecinas, saltos grandes de pesos entre ellas
suben la energ\'ia y son penalizados.

Energ\'ia de suavidad de una fila de pesos:
\[
\mathcal E_i = \bm w_i^T\K_G\bm w_i.
\]
Total:
\[
\mathcal S_G(\W)=\sum_i \bm w_i^T\K_G\bm w_i
=\operatorname{tr}(\W\K_G\W^T).
\]

\subsection{Objetivo incremental regularizado}
Para una poda tentativa:
\[
\Delta\W=\W_{\text{new}}-\W_{\text{old}},
\]
se minimiza:
\[
\frac{1}{2}\|\Delta\W\|_F^2
\;+\;
\frac{\alpha_G}{2}\operatorname{tr}(\Delta\W \K_G \Delta\W^T),
\]
sujeto a exactitud local + clamps + positividad final.

Vectorizando $\bm z=\operatorname{vec}(\W_{\text{new}})$:
\[
\bm H=
\bm I+\alpha_G(\K_G\otimes \bm I_r),
\]
\[
\min_{\bm z}\;\frac12 \bm z^T\bm H\bm z-\bm g^T\bm z,
\quad \bm g=\bm H\bm z_{\text{old}}
\;(\text{modo incremental}).
\]

Aqu\'i:
\begin{itemize}[leftmargin=0.8cm]
    \item $\operatorname{vec}(\W)$ apila columnas de $\W$ en un solo vector.
    \item $r$ = n\'umero de candidatos activos en esa iteraci\'on de pruning.
    \item $\bm I_r$ = identidad $r\times r$.
    \item $\K_G\otimes \bm I_r$ = acopla nodos latentes, pero sin mezclar ``identidad de elemento''.
\end{itemize}

Ejemplo r\'apido de \texttt{vec}: si
\[
\W=
\begin{bmatrix}
1 & 3\\
2 & 4
\end{bmatrix},
\]
entonces
\[
\operatorname{vec}(\W)=
\begin{bmatrix}
1\\2\\3\\4
\end{bmatrix}.
\]

\subsection{Parametrizaci\'on en nullspace (Option B)}
Con active sets fijos en una iteraci\'on:
\[
\bm z = \bm z_p + \bm N \bm y,
\]
y se resuelve:
\[
(\bm N^T\bm H\bm N)\bm y = \bm N^T(\bm g-\bm H\bm z_p).
\]
Luego se reconstruye $\W_{\text{new}}$, se revisan negativos y se actualizan
$\mathcal D_j$ nodo a nodo.

\subsection{Criterios de ranking de candidatos (importante)}
En \texttt{GraphBased\_MAW\_ECM\_2v.m} aparece:
\begin{itemize}[leftmargin=0.8cm]
    \item criterio 0: factibilidad/clamps (legado),
    \item criterio 1: energ\'ia suave $S$,
    \item criterio 2: robustez anti-picos con percentil 95 de saltos ($R95$) y
    desempate con $R_{\max}$.
\end{itemize}

Con aristas $(a,b)$ del grafo y
$\Delta\W=\W_{\text{trial}}-\W_{\text{old}}$:
\[
J_{i,e}=\Delta W_{i,a}-\Delta W_{i,b},
\quad
R95=\operatorname{percentil}_{95}(|J_{i,e}|),
\quad
R_{\max}=\max |J_{i,e}|.
\]
Esto busca evitar soluciones tipo ``aguja'' (\emph{needle}).

\section{Regresi\'on final: de muestras a funci\'on de pesos}
Script: \texttt{RegressionViaRBF\_master\_weights.m}.

Al final del pruning se tiene $\W_{\text{red}}$ en muestras de entrenamiento
$\{\bm q_M^{(j)}\}$.
Se ajusta un RBF gaussiano anisotr\'opico por peso:
\[
w_i(\bm q_M)\approx
\sum_{k=1}^{n_{\text{RBF}}}\gamma_{ik}\,\phi_k(\bm q_M)+p_i(\bm q_M),
\]
\[
\phi_k(\bm q_M)=
\exp\!\left(
-\frac12(\bm q_M-\bm c_k)^T\D_k^{-2}(\bm q_M-\bm c_k)
\right).
\]

En forma matricial:
\[
\begin{bmatrix}\Phi & P\end{bmatrix}
\begin{bmatrix}\bm\gamma_i\\ \bm\eta_i\end{bmatrix}
\approx \bm w_i.
\]

Desglose expl\'icito:
\begin{itemize}[leftmargin=0.8cm]
    \item $\bm w_i\in\mathbb R^{N_M}$ = valores del peso del elemento $i$ en las $N_M$ muestras latentes.
    \item $\Phi\in\mathbb R^{N_M\times n_{\text{RBF}}}$ = matriz de kernels RBF evaluados en muestras/centros.
    \item $P$ = base polin\'omica (constante o af\'in, seg\'un opci\'on).
    \item $\bm\gamma_i,\bm\eta_i$ = coeficientes a ajustar para el peso del elemento $i$.
\end{itemize}

El c\'odigo adem\'as normaliza por RMS cada salida para mejorar condici\'on
num\'erica.

\section{Ejemplo did\'actico 2D peque\~no}
Sup\'on:
\begin{itemize}[leftmargin=0.8cm]
    \item malla latente $3\times 3$ $\Rightarrow N_M=9$,
    \item candidatos iniciales $n_c=12$ (de ECM fijo),
    \item restricciones locales de tama\~no $r_j=5$ (aprox. constante).
\end{itemize}

\textbf{Inicio:}
\[
\W^{(0)}=\bm w_0\,\bm 1_{1\times 9},
\]
misma columna en los 9 nodos latentes.

\textbf{Prueba de poda de un candidato $p$:}
\begin{enumerate}[leftmargin=1.2cm]
    \item fijar fila $p$ a cero en todos los nodos,
    \item resolver redistribuci\'on (Etapa 1),
    \item si hay negativos, pasar a Etapa 2 o 3,
    \item si factible en todos los nodos, aceptar poda.
\end{enumerate}

\textbf{Interpretaci\'on geom\'etrica:}
la poda es global (fila completa), pero la reparaci\'on por positividad es
local (entry por nodo), y la Etapa 3 a\~nade acoplamiento entre nodos vecinos.

\section{Conexi\'on directa con implementaci\'on de Stage 12 (residual-only)}
En nuestro flujo residual-only, la lectura pr\'actica de las matrices es:
\[
\A^{(s)}_{\text{res}} \in \mathbb R^{n_q\times n_{\text{el}}},
\qquad
\bm w^{(s)}\in\mathbb R^{n_{\text{el}}},
\qquad
\bm b^{(s)}_{\text{res}}\in\mathbb R^{n_q},
\]
y la ecuaci\'on objetivo es
\[
\A^{(s)}_{\text{res}}\bm w^{(s)} \approx \bm b^{(s)}_{\text{res}}.
\]
Aqu\'i:
\begin{itemize}[leftmargin=0.8cm]
    \item $s$ = snapshot (paso de entrenamiento),
    \item $n_q$ = dimensi\'on de coordenadas reducidas primarias ($q_p$),
    \item $n_{\text{el}}$ = n\'umero de elementos candidatos del soporte actual,
    \item fila/columna de $\A_{\text{res}}^{(s)}$: fila = componente de residual proyectado; columna = elemento.
\end{itemize}
Si lo miras en los archivos de dataset, esto es justamente el bloque por snapshot de
\texttt{Q\_ecm.dat} y su objetivo correspondiente en \texttt{b\_full.dat}.

Si en tu pipeline actual quieres usar s\'olo residual:
\begin{enumerate}[leftmargin=1.2cm]
    \item construir datasets de residual ($Q_{\text{ecm}},b$),
    \item ejecutar MAW--ECM en residual para obtener soporte y pesos adaptativos,
    \item ajustar RBF para $w_{\text{res}}(\bm q_p)$,
    \item online: evaluar s\'olo residual hiperrreducido y comparar $q_p$
    PROM-GPR vs HPROM-GPR.
\end{enumerate}

Esto es consistente con la l\'ogica del profesor: primero asegurar que
\textbf{el mecanismo de soporte/pesos} funciona en un objetivo, luego extender.

\section{Qu\'e debes mirar siempre para diagnosticar}
\begin{enumerate}[leftmargin=1.2cm]
    \item \textbf{Factibilidad local}: cu\'antos candidatos se rechazan por no
    poder satisfacer exactitud + positividad.
    \item \textbf{Evoluci\'on de active sets}: si hay mucho clamp local, esperar
    no suavidad.
    \item \textbf{M\'etricas de roughness} ($S$, $R95$, $R_{\max}$): indican si
    aparecen agujas.
    \item \textbf{Calidad de regresi\'on RBF}: error train y estabilidad al
    evaluar puntos intermedios.
    \item \textbf{Compresi\'on real}: $|\mathcal Z_{\text{red}}|$ frente a
    error en QoI (aqu\'i, trayectoria $q_p$).
\end{enumerate}

\section{Resumen final corto}
MAW--ECM de tu profesor no es ``ECM + regresi\'on''.
Es un \textbf{proceso de poda factible con restricciones}:
\begin{itemize}[leftmargin=0.8cm]
    \item soporte com\'un global,
    \item pesos adaptativos por estado latente,
    \item exactitud local en cada estado,
    \item positividad (active-set),
    \item suavidad espacial en el manifold (grafo),
    \item y, al final, una regresi\'on RBF para uso online.
\end{itemize}

Esa es la l\'ogica de fondo que debes conservar al pasar de 2D a 3D: no copiar
una forma superficial, sino mantener esta estructura matem\'atica y algor\'itmica.

\end{document}
